\scanpage{63}
\mspage{19}
La question essentielle est donc comment \unsure{reconstituer|réconstituer}, pour un graphe \deleted{$\unsure{\Gamma|\mathrm{top.}}$} pondéré :
\deleted{à repères \unsure{donné|donnés}} \inserted{\unsure{repères|repéré}} (disons) $\underline\Gamma$, le
groupoïde $\mathcal T_{\mathrm{top}}^\pm(\underline\Gamma)$,
\deleted{plus} \inserted{avec} sa dépendance fonctorielle en $\underline\Gamma$,
\deleted{\unsure{des foncteurs les|des foncteurs des}} \inserted{et l'explicitation}
des sept opérations, à l'exception de $4^\circ)$ et $5^\circ)$.
\margin{\uline{NB} on laisse ici tomber le ``$S$'' ou ``$D$'', \unsure{contenus|contenu} déjà dans $\underline\Gamma$}

Pour l'opération $4^\circ)$, elle sera obtenue automatiquement, par la description explicite des
$\mathcal T_{\mathrm{top}}^\pm(\underline\Gamma)$
(qui sera faite en $=$ temps et de la $=$ façon que pour les ${}^{\unsure{S}}\mathcal T_{g,\nu}^\pm$ et $S\mathcal T_{g,\nu}^\pm$ \unsure{connus|comme ens|ci-dessus}),
et le bouchage des trous \inserted{pourrait} s'expliciter à l'aide \deleted{\unsure{de ces|des} opérations} \inserted{des cas} particuliers de ces opérations pour les
\deleted{\unsure{catégories|analogues}\unsure{,} ou}
$S\mathcal T_{g,\nu}$ et $\mathcal T_{g,\nu}$, qu'il faudrait jeter dans le sac des
\deleted{\unsure{fl.} \unsure{(}des \unsure{unités|données}} \inserted{structures} ``élémentaires'' : diminution

\diagram{\begin{tikzcd}[ampersand replacement=\&,column sep=huge,
execute at end picture={
\node[anchor=east,inner sep=0pt] at ([xshift=-5mm]\tikzcdmatrixname-1-1.west) {$\unsure{(**)|\deleted{(**)}}$};
\draw ([xshift=2mm,yshift=-1.5mm]\tikzcdmatrixname-1-1.south west)
 -- ([xshift=2mm,yshift=-3mm]\tikzcdmatrixname-1-1.south west)
 -- ([yshift=-3mm]\tikzcdmatrixname-1-1.south east)
 -- ([yshift=-1.5mm]\tikzcdmatrixname-1-1.south east);
\node[anchor=north west,inner sep=0pt,align=left] at ([xshift=2mm,yshift=-5mm]\tikzcdmatrixname-1-1.south west) {\inserted{\begin{tabular}{@{}l@{}}
\unsure{les|un} $X\in\operatorname{Ob}\mathcal T_{g,\nu}^\pm$, \unsure{avec|ou}\\
\unsure{\deleted{la}|\deleted{les}|la} \unsure{donnée|donné} d'un\\
sous-ens. $I'\subset I(X)$\\
avec $\operatorname{card}I'=\nu'$
\end{tabular}}};
}]
\deleted{S}\mathcal T_{g,\nu,\nu'}^\pm\arrow[r]
\&\deleted{S}\mathcal T_{g,\nu\unsure{\pm|-|+}\nu'}
\end{tikzcd}}
\margin{\figure{\begin{tikzpicture}[x=1cm,y=-1cm]
\draw[line width=.45pt] (.18,.05)--(.20,.74)--(.13,1.69)--(.08,2.35)--(.10,3.20)--(.04,4.37)--(.05,5.31) .. controls (.06,5.70) and (.02,5.83) .. (-.08,5.92);
\node[anchor=north west,align=left] at (.50,.36) {avec ou\\sans $S$ ?};
\end{tikzpicture}}}
\mspage{62}
